% GENERATED FILE. Edit canonical lessons, then run npm run generate:books.
% canonical-source-hash: fnv1a64:72387d3ebfd8034d
% canonical-lessons: MW-C290-bahi, MW-C290-rojgar, MW-C290-saru, MW-C290-khatam, MW-C290-abhyas, MW-R290-review-285-290, MW-R290-review-final-290

\chapter{Starting and Finishing}
\label{ch:mw-starting-and-finishing}
\hlchaptermodality{290}

\begin{chapteropening}
\textbf{By the end of this chapter:} \emph{I can say when work starts and when it is finished.}

Complete the last lesson of chapter 290: I can say when work starts and when it is finished.
\end{chapteropening}

\section[\texorpdfstring{\mw{बही} (bahī)}{बही (bahī)}]{\mw{बही} (bahī) — a ledger}

\label{lesson:MW-C290-bahi}



\begin{quote}
Before the new one: say the Marwadi for a carpenter's plane, then the Marwadi for to paint a picture.
\end{quote}

\subsection*{You'll want to know: \mw{बही}}
\textbf{\mw{बही}} — \emph{bahī} — \textquotedblleft{}a ledger\textquotedblright{}.

\begin{quote}
\mw{मालक} \mw{बही} \mw{राखै} \mw{है।}
\end{quote}

\begin{quote}
\emph{mālak bahī rākhai hai.}
\end{quote}

\begin{quote}
The owner keeps a ledger.
\end{quote}

A feminine noun.

\begin{cousinweb}
\textbf{\mw{बही}} (\emph{bahī}) is inherited from Sanskrit \emph{vahikā}, found in a word for a king's diary; CDIAL gives no Marwari form in this entry, but lists Hindi \emph{bahī}, \textquotedblleft{}account-book, ledger\textquotedblright{} (CDIAL 11460).
\end{cousinweb}

\begin{grammarlens}[title={What you've built}]
One more word for this chapter.
\end{grammarlens}

\subsection*{Your turn}
\begin{itemize}
\raggedright
  \item \emph{Say it:} \emph{bahī}
  \item \emph{Say it:} \emph{bahī}, once more
  \item \emph{Say it:} \emph{mālak bahī rākhai hai.} — The owner keeps a ledger.
  \item \emph{Recall:} say the Marwadi for a carpenter's plane, then the Marwadi for to paint a picture, then say \emph{bahī} again
\end{itemize}

\begin{tcolorbox}[breakable,colback=teal!4,colframe=teal!35!black,title={Before you move on}]
What does \mw{बही} mean? (\textquotedblleft{}A ledger\textquotedblright{}.) Say it once more.
\end{tcolorbox}
\section[\texorpdfstring{\mw{रोजगार} (rojgār)}{रोजगार (rojgār)}]{\mw{रोजगार} (rojgār) — employment}

\label{lesson:MW-C290-rojgar}



\begin{quote}
Before the new one: say the Marwadi for to paint a picture, then the Marwadi for a ledger.
\end{quote}

\subsection*{You'll want to know: \mw{रोजगार}}
\textbf{\mw{रोजगार}} — \emph{rojgār} — \textquotedblleft{}employment\textquotedblright{}.

\begin{quote}
\mw{सहर} \mw{मांय} \mw{रोजगार} \mw{है।}
\end{quote}

\begin{quote}
\emph{sahar mā̃y rojgār hai.}
\end{quote}

\begin{quote}
There is work to be had in the city.
\end{quote}

A masculine noun.

\begin{cousinweb}
\textbf{\mw{रोजगार}} (\emph{rojgār}) is usually traced to Persian \emph{rōzgār}, \textquotedblleft{}time; livelihood\textquotedblright{}.
\end{cousinweb}

\begin{grammarlens}[title={What you've built}]
One more word for this chapter.
\end{grammarlens}

\subsection*{Your turn}
\begin{itemize}
\raggedright
  \item \emph{Say it:} \emph{rojgār}
  \item \emph{Say it:} \emph{rojgār}, once more
  \item \emph{Say it:} \emph{sahar mā̃y rojgār hai.} — There is work to be had in the city.
  \item \emph{Recall:} say the Marwadi for to paint a picture, then the Marwadi for a ledger, then say \emph{rojgār} again
\end{itemize}

\begin{tcolorbox}[breakable,colback=teal!4,colframe=teal!35!black,title={Before you move on}]
What does \mw{रोजगार} mean? (\textquotedblleft{}Employment\textquotedblright{}.) Say it once more.
\end{tcolorbox}
\section[\texorpdfstring{\mw{सरू} (sarū)}{सरू (sarū)}]{\mw{सरू} (sarū) — begun; the start}

\label{lesson:MW-C290-saru}



\begin{quote}
Before the new one: say the Marwadi for a ledger, then the Marwadi for employment.
\end{quote}

\subsection*{You'll want to know: \mw{सरू}}
\textbf{\mw{सरू}} — \emph{sarū} — \textquotedblleft{}begun; the start\textquotedblright{}.

\begin{quote}
\mw{पाठ} \mw{सरू} \mw{करो।}
\end{quote}

\begin{quote}
\emph{pāṭh sarū karo.}
\end{quote}

\begin{quote}
Start the lesson.
\end{quote}

\begin{cousinweb}
\textbf{\mw{सरू}} (\emph{sarū}) is usually traced, through Persian \emph{shurūʿ}, to Arabic \emph{shurūʿ}, \textquotedblleft{}a beginning\textquotedblright{}. With \textbf{\mw{करो}} it asks someone to start.
\end{cousinweb}

\begin{grammarlens}[title={What you've built}]
One more word for this chapter.
\end{grammarlens}

\subsection*{Your turn}
\begin{itemize}
\raggedright
  \item \emph{Say it:} \emph{sarū}
  \item \emph{Say it:} \emph{sarū}, once more
  \item \emph{Say it:} \emph{pāṭh sarū karo.} — Start the lesson.
  \item \emph{Recall:} say the Marwadi for a ledger, then the Marwadi for employment, then say \emph{sarū} again
\end{itemize}

\begin{tcolorbox}[breakable,colback=teal!4,colframe=teal!35!black,title={Before you move on}]
What does \mw{सरू} mean? (\textquotedblleft{}Begun; the start\textquotedblright{}.) Say it once more.
\end{tcolorbox}
\section[\texorpdfstring{\mw{खतम} (khatam)}{खतम (khatam)}]{\mw{खतम} (khatam) — finished}

\label{lesson:MW-C290-khatam}



\begin{quote}
Before the new one: say the Marwadi for employment, then the Marwadi for begun.
\end{quote}

\subsection*{You'll want to know: \mw{खतम}}
\textbf{\mw{खतम}} — \emph{khatam} — \textquotedblleft{}finished\textquotedblright{}.

\begin{quote}
\mw{म्हारो} \mw{पाठ} \mw{खतम} \mw{है।}
\end{quote}

\begin{quote}
\emph{mhāro pāṭh khatam hai.}
\end{quote}

\begin{quote}
My lesson is finished.
\end{quote}

\begin{cousinweb}
\textbf{\mw{खतम}} (\emph{khatam}) is usually traced, through Persian \emph{khatm}, to Arabic \emph{khatm}, \textquotedblleft{}sealing; the end\textquotedblright{}.
\end{cousinweb}

\begin{grammarlens}[title={What you've built}]
One more word for this chapter.
\end{grammarlens}

\subsection*{Your turn}
\begin{itemize}
\raggedright
  \item \emph{Say it:} \emph{khatam}
  \item \emph{Say it:} \emph{khatam}, once more
  \item \emph{Say it:} \emph{mhāro pāṭh khatam hai.} — My lesson is finished.
  \item \emph{Recall:} say the Marwadi for employment, then the Marwadi for begun, then say \emph{khatam} again
\end{itemize}

\begin{tcolorbox}[breakable,colback=teal!4,colframe=teal!35!black,title={Before you move on}]
What does \mw{खतम} mean? (\textquotedblleft{}Finished\textquotedblright{}.) Say it once more.
\end{tcolorbox}
\section[\texorpdfstring{\mw{अभ्यास} (abhyās)}{अभ्यास (abhyās)}]{\mw{अभ्यास} (abhyās) — practice}

\label{lesson:MW-C290-abhyas}



\begin{quote}
Before the new one: say the Marwadi for begun, then the Marwadi for finished.
\end{quote}

\subsection*{You'll want to know: \mw{अभ्यास}}
\textbf{\mw{अभ्यास}} — \emph{abhyās} — \textquotedblleft{}practice\textquotedblright{}.

\begin{quote}
\mw{रोज} \mw{अभ्यास} \mw{करो।}
\end{quote}

\begin{quote}
\emph{roj abhyās karo.}
\end{quote}

\begin{quote}
Practise every day.
\end{quote}

A masculine noun.

\begin{cousinweb}
\textbf{\mw{अभ्यास}} (\emph{abhyās}) is a tatsama, taken from Sanskrit \emph{abhyāsa}, \textquotedblleft{}repetition; repeated exercise\textquotedblright{} (Monier-Williams).
\end{cousinweb}

\begin{grammarlens}[title={What you've built}]
One more word for this chapter.
\end{grammarlens}

\subsection*{Your turn}
\begin{itemize}
\raggedright
  \item \emph{Say it:} \emph{abhyās}
  \item \emph{Say it:} \emph{abhyās}, once more
  \item \emph{Say it:} \emph{roj abhyās karo.} — Practise every day.
  \item \emph{Recall:} say the Marwadi for begun, then the Marwadi for finished, then say \emph{abhyās} again
\end{itemize}

\begin{tcolorbox}[breakable,colback=teal!4,colframe=teal!35!black,title={Before you move on}]
What does \mw{अभ्यास} mean? (\textquotedblleft{}Practice\textquotedblright{}.) Say it once more.
\end{tcolorbox}
\section[(dialogue)]{Looking back: from Learning by Heart to Starting and Finishing}

\label{lesson:MW-R290-review-285-290}



\begin{quote}
Say the Marwadi for a letter of the alphabet and a page.
\end{quote}

\subsection*{Your turn}
Say these aloud:

\begin{itemize}
\raggedright
  \item a letter of the alphabet, a page, a lesson, to learn by heart, to pronounce
  \item an exam, sums, a result, a class at school, to mark schoolwork
  \item to teach a skill, to teach lessons, a guru, a pupil, a wrestling ground
  \item hard work, capable, a master of a craft, a trade, a labourer
  \item a machine, to drive, to rub smooth, a carpenter's plane, to paint a picture
  \item a ledger, employment, begun, finished, practice
\end{itemize}

\begin{tcolorbox}[breakable,colback=teal!4,colframe=teal!35!black,title={Before you move on}]
A letter of the alphabet? (\textbf{\mw{आखर}}, \emph{ākhar}.) A page? (\textbf{\mw{पानो}}, \emph{pāno}.) A lesson? (\textbf{\mw{पाठ}}, \emph{pāṭh}.) To learn by heart? (\textbf{\mw{रटणो}}, \emph{raṭṇo}.) To pronounce? (\textbf{\mw{उचारणो}}, \emph{uchārṇo}.) An exam? (\textbf{\mw{इम्तिहान}}, \emph{imtihān}.) Sums? (\textbf{\mw{हिसाब}}, \emph{hisāb}.) A result? (\textbf{\mw{नतीजो}}, \emph{natījo}.) A class at school? (\textbf{\mw{दरजो}}, \emph{darjo}.) To mark schoolwork? (\textbf{\mw{आंकणो}}, \emph{ā̃kṇo}.) To teach a skill? (\textbf{\mw{सिखावणो}}, \emph{sikhāvṇo}.) To teach lessons? (\textbf{\mw{भणावणो}}, \emph{bhaṇāvṇo}.) A guru? (\textbf{\mw{गुरु}}, \emph{guru}.) A pupil? (\textbf{\mw{चेलो}}, \emph{chelo}.) A wrestling ground? (\textbf{\mw{अखाड़ो}}, \emph{akhāṛo}.) Hard work? (\textbf{\mw{मेहनत}}, \emph{mehnat}.) Capable? (\textbf{\mw{लायक}}, \emph{lāyak}.) A master of a craft? (\textbf{\mw{उस्ताद}}, \emph{ustād}.) A trade? (\textbf{\mw{धंधो}}, \emph{dhandho}.) A labourer? (\textbf{\mw{मजूर}}, \emph{majūr}.) A machine? (\textbf{\mw{मसीन}}, \emph{masīn}.) To drive? (\textbf{\mw{चलावणो}}, \emph{chalāvṇo}.) To rub smooth? (\textbf{\mw{घिसणो}}, \emph{ghisṇo}.) A carpenter's plane? (\textbf{\mw{रंदो}}, \emph{rando}.) To paint a picture? (\textbf{\mw{चीतरणो}}, \emph{chītarṇo}.) A ledger? (\textbf{\mw{बही}}, \emph{bahī}.) Employment? (\textbf{\mw{रोजगार}}, \emph{rojgār}.) Begun? (\textbf{\mw{सरू}}, \emph{sarū}.) Finished? (\textbf{\mw{खतम}}, \emph{khatam}.) Practice? (\textbf{\mw{अभ्यास}}, \emph{abhyās}.)
\end{tcolorbox}
\section[(dialogue)]{Starting and Finishing — once more}

\label{lesson:MW-R290-review-final-290}



\begin{quote}
Say the Marwadi for a ledger and employment.
\end{quote}

\subsection*{Your turn}
\begin{itemize}
\raggedright
  \item \emph{Say it:} a ledger, employment, begun, finished, practice
\end{itemize}

\begin{tcolorbox}[breakable,colback=teal!4,colframe=teal!35!black,title={Before you move on}]
A ledger? (\textbf{\mw{बही}}, \emph{bahī}.) Employment? (\textbf{\mw{रोजगार}}, \emph{rojgār}.) Begun? (\textbf{\mw{सरू}}, \emph{sarū}.) Finished? (\textbf{\mw{खतम}}, \emph{khatam}.) Practice? (\textbf{\mw{अभ्यास}}, \emph{abhyās}.)
\end{tcolorbox}
